\documentclass{situationist-preprint}

\title{Twin Earth Revisited}
\msid{LNG-08-b}
\msfamily{LNG-08}
\msdiscipline{Philosophy of Language}
\mstype{Companion manuscript}
\mskeywords{Twin Earth; externalism; natural kind terms; Putnam; Burge; reference-fixing profile}

\begin{document}
\maketitle

\begin{abstract}
LNG-08-a argues that meanings are not in the head but depend on it: the meaning of a kind term is fixed jointly by the environment, the community and a reference-fixing profile in the speaker's internal state. Putnam's original Twin Earth case establishes only the first half. We run five variations of the case, each holding fixed some of the factors and varying the others: a social variation, a sortal variation, a mixed environment, a deferential variation and a jade-type case. In each, the verdict that intuition and scientific practice deliver is the verdict the joint account predicts, and in none is it the verdict of either a purely internalist or a purely environmental account.
\end{abstract}

\section{The base case}

Oscar on Earth and his internal duplicate on Twin Earth, in 1750, both say \enquote{water}; Oscar's word applies to H$_2$O, his twin's to XYZ \citep{Putnam1975}. This is Principle~P1 of LNG-08-a, from which Theorem~3.1 there infers that meanings are not in the head. Proposition~3.3 there states the positive account: speakers who agree in internal state, environment and community mean the same by a kind term. The internal state contributes the reference-fixing profile of Definition~2.4: the sortal under which the samples are taken, the deferential and sameness intentions, and the associated information. We test the account by varying one factor at a time.

\section{Five variations}

\emph{Social variation.} Burge's patient believes that he has arthritis in his thigh. In the actual community \enquote{arthritis} applies only to inflammation of the joints; in a counterfactual community it applies to rheumatoid ailments generally \citep{Burge1979}. The patient's internal state and physical environment are fixed; the community varies; his meaning varies with it. The joint account predicts this through his deferential intention, clause (b) of Definition~2.4: he intends to use the word as those from whom he acquired it do, and they use it differently in the two communities. A purely environmental account, with the community left out, predicts no difference.

\emph{Sortal variation.} Two speakers in the same environment ground a term in the same samples of water, one taking them as samples of a kind of substance, the other as samples of a kind of liquid that people drink from springs. Asked later about ice, or about distilled water from a laboratory, the first applies the term and the second does not. Principle~P2 of LNG-08-a records the verdict, and Theorem~3.2 there infers that meaning depends on the head: environment and community are the same, the sortal differs, and the extension differs. No purely environmental account can deliver this verdict.

\emph{Mixed environment.} Suppose the lakes of Earth contained H$_2$O and XYZ in roughly equal measure, and the term was grounded in samples of both. The sameness relation appropriate to the sortal \enquote{substance} relates a sample to what shares its chemical nature; the samples share none. The causal component of the practice then fails to single out one kind, as in clause (b) of Definition~3.3 of LNG-02-a. On the hybrid theory of LNG-02-a, the descriptive component takes over, and the associated information (clause (c) of Definition~2.4 of LNG-08-a), which describes a clear, drinkable liquid filling lakes, applies to both. The word applies to both kinds. This agrees with Putnam's own verdict on jade, discussed below, and shows the associated information doing work exactly where the grounding fails, as LNG-08-a predicts.

\emph{Deferential variation.} A speaker introduces \enquote{water} into her idiolect by stipulation, intending to apply it to whatever is clear, drinkable and fills lakes, and explicitly not to defer to any community or to the underlying nature of samples. On Twin Earth her duplicate does the same. Since neither has a sameness intention tied to a sortal, nor a deferential intention, the joint account predicts that they mean the same: their terms apply to whatever satisfies the stipulated description, H$_2$O and XYZ alike. The verdict seems right. It shows that the environment fixes extension only through the intentions in the profile; remove those intentions and the environmental difference makes no difference.

\emph{Jade.} \enquote{Jade} was applied to samples of two minerals, jadeite and nephrite, before their difference was known, and Putnam held that it applies to both \citep{Putnam1975}. This is the mixed environment realized. The joint account predicts the verdict for the same reason; a purely internalist account predicts it too, but fails on the base case.

\section{What the variations show}

Each variation keeps two of the three factors fixed and varies the third, or varies a clause of the profile. The internal state is shown to matter by the sortal and deferential variations, the community by the social variation, the environment by the base case. The mixed environment and jade show that the associated information becomes relevant when, and only when, the causal grounding fails to single out a kind, which is why kind terms refer independently of the descriptions associated with them in the normal case \citep{Kripke1980}.

\begin{objection}
Narrow content theorists hold that the twins share a content that is in the head, and that the variations show only that wide content varies \citep{Fodor1987}.
\end{objection}

\begin{reply}
The variations are compatible with there being a narrow state shared by the twins: the profile of Definition~2.4 of LNG-08-a is such a state. What they deny is that this state is a meaning, in the sense of Definition~2.1 there, since it does not determine extension. Calling it a content does not change what fixes extension.
\end{reply}

\begin{objection}
The deferential variation is artificial; ordinary speakers always defer.
\end{objection}

\begin{reply}
Ordinary speakers usually defer, which is why environmental and social factors usually matter. The variation shows that their mattering runs through an internal state, the deferential intention, that could have been absent. That is what the claim that meanings depend on the head requires.
\end{reply}

\section{Conclusion}

In each variation of Twin Earth, the meaning of the kind term is the one fixed jointly by environment, community and the speaker's reference-fixing profile. Meanings are not in the head, and they depend on it.

\begin{thebibliography}{9}\small
\bibitem[Burge(1979)]{Burge1979} Burge, T. (1979). Individualism and the mental. \emph{Midwest Studies in Philosophy}, 4, 73--121.
\bibitem[Fodor(1987)]{Fodor1987} Fodor, J.~A. (1987). \emph{Psychosemantics}. MIT Press.
\bibitem[Kripke(1980)]{Kripke1980} Kripke, S.~A. (1980). \emph{Naming and Necessity}. Harvard University Press.
\bibitem[Putnam(1975)]{Putnam1975} Putnam, H. (1975). The meaning of \enquote{meaning}. In K.~Gunderson (ed.), \emph{Language, Mind, and Knowledge}, Minnesota Studies in the Philosophy of Science, Vol.~7. University of Minnesota Press.
\end{thebibliography}

\end{document}
